\chapter{Related Works}

\section{API-First and Contract-Driven Engineering}

API-first engineering treats the API specification as the primary artifact for implementation, testing, and documentation~\cite{openapi_spec}. OpenAPI/Swagger has become the standard contract format for machine-readable endpoint and schema definition.

In contract-driven practice, producer and consumer implementations must remain aligned with that specification. This is critical in enterprise integration, where REST APIs often front legacy EJB-based services. The platform in this thesis follows this model: it parses OpenAPI/Swagger contracts, derives structured metadata, and drives mapping and generation from that contract as the source of truth.

\section{Existing Code Generation Ecosystems}

OpenAPI Generator and Swagger Codegen show that specification-driven scaffolding can accelerate client and server development. They are effective for broad language coverage and template customization.

However, these generators are optimized for generic or greenfield outputs. Enterprise connector environments impose fixed module hierarchies, legacy EJB service contracts, and organization-specific Maven pipelines. Template-based and model-driven approaches remain useful, but they often require substantial adaptation before they fit this context. The examined platform therefore applies a direct contract-to-enterprise transformation flow that combines parsing, metadata extraction, mapping, and generation in one pipeline.

\section{Bytecode Introspection and Metadata Extraction}

The platform's key differentiator is bytecode introspection of compiled client-kit JAR files without source code. Using \texttt{java-class-tools}, it parses class metadata in Node.js/browser environments and avoids JVM-only tooling.

Its extraction depth goes beyond signatures: it resolves hierarchy information, detects EJB interfaces through annotation and naming heuristics, parses generic signatures, flattens nested field paths, and normalizes synthetic parameter names. This output directly supports accurate interactive mapping.

\section{Workflow Automation for Integration Teams}

CI/CD systems automate downstream build and deployment, while low-code integration platforms automate high-level workflow design. Neither category directly addresses the full upstream connector-preparation workflow in legacy EJB environments.

The studied platform sits between these approaches. It offers guided mapping and automation comparable to low-code usability while preserving control required for organization-specific module structures and build chains.

\section{Real-Time Communication in Development Tools}

Real-time feedback is important for long-running generation tasks. Socket.IO~\cite{socketio_docs} provides event messaging, reconnection, and fallback behaviour that fits enterprise network conditions.

In this platform, Socket.IO is used for (i) line-by-line generation log streaming and (ii) state synchronisation events for modules, specifications, client kits, and configurations, with polling fallback for resilience.

\section{Comparison and Positioning}

Existing tools cover fragments of the problem: generic generators focus on code templates, low-code tools focus on abstraction, and CI/CD tools focus on downstream automation. The studied platform combines contract parsing, bytecode extraction, interactive mapping, YAML generation, and scripted module creation in one enterprise-conformant workflow.

Its position is therefore distinct: compiled enterprise artifacts are first-class inputs, outputs match existing module conventions, and users receive real-time operational visibility.

Compared with generic generators, low-code platforms, and CI/CD pipelines, the platform studied here is distinctive in three combined capabilities: bytecode-driven enterprise metadata extraction, interactive contract-to-EJB mapping, and generation output aligned to existing module conventions in air-gapped deployment settings.
